\documentclass[12pt,a4paper]{article}

\usepackage[left=30mm,top=20mm,right=15mm,bottom=33mm,includeheadfoot]{geometry}
\usepackage[utf8]{inputenc}
\usepackage[T2A]{fontenc}
\usepackage[russian,english]{babel}
\usepackage{setspace}
\usepackage{indentfirst}
\usepackage{array}
\usepackage{tabularx}
\usepackage{longtable}
\usepackage{multirow}
\usepackage{makecell}
\usepackage{graphicx}
\usepackage{hyperref}
\usepackage{xurl}
\usepackage{fancyhdr}
\usepackage{titlesec}
\usepackage{enumitem}
\usepackage{lastpage}
\usepackage[absolute,overlay]{textpos}

\hypersetup{hidelinks}
\onehalfspacing
\setlength{\parindent}{1.25cm}
\setlength{\headheight}{30pt}
\setlength{\footskip}{72pt}
\emergencystretch=3em
\setlength{\TPHorizModule}{1mm}
\setlength{\TPVertModule}{1mm}
\textblockorigin{0mm}{0mm}

\newcommand{\doccodebase}{RU.17701729.05.05-01}
\newcommand{\doccode}{RU.17701729.05.05-01 ТЗ 01-1}
\newcommand{\doccodeapproval}{RU.17701729.05.05-01 ТЗ 01-1-ЛУ}
\newcommand{\doctitle}{Андроид-приложение для совместного планирования поездок}

\titleformat{\section}{\large\bfseries\centering}{\thesection}{1em}{\MakeUppercase}
\titleformat{\subsection}{\normalsize\bfseries}{\thesubsection}{1em}{}
\titleformat{\subsubsection}{\normalsize\bfseries}{\thesubsubsection}{1em}{}
\titlespacing*{\section}{0pt}{1.5ex plus 0.5ex minus 0.3ex}{1.0ex}
\titlespacing*{\subsection}{0pt}{1.2ex plus 0.4ex minus 0.2ex}{0.8ex}
\titlespacing*{\subsubsection}{0pt}{1.0ex plus 0.3ex minus 0.2ex}{0.5ex}

\setlist[enumerate]{leftmargin=1.2cm,itemsep=0.2em,topsep=0.3em}
\setlist[itemize]{leftmargin=1.2cm,itemsep=0.2em,topsep=0.3em}

\newcommand{\pagefooter}{%
  \begingroup
  \setlength{\tabcolsep}{3pt}%
  \renewcommand{\arraystretch}{1.0}%
  \footnotesize
  \begin{tabularx}{\textwidth}{|>{\centering\arraybackslash}p{0.305\textwidth}|>{\centering\arraybackslash}p{0.14\textwidth}|>{\centering\arraybackslash}p{0.13\textwidth}|>{\centering\arraybackslash}p{0.125\textwidth}|>{\centering\arraybackslash}X|}
    \hline
    Изм. & Лист & № докум. & Подп. & Дата \\
    \hline
    \doccode & & & & \\
    \hline
    Инв.~№~подл. & Подп.~и~дата & Взам.~инв.~№ & Инв.~№~дубл. & Подп.~и~дата \\
    \hline
  \end{tabularx}
  \endgroup
}

\fancypagestyle{tzstyle}{
  \fancyhf{}
  \fancyhead[C]{\textbf{\thepage}\\[-0.2ex]\textbf{\doccode}}
  \fancyfoot[C]{\pagefooter}
  \renewcommand{\headrulewidth}{0pt}
  \renewcommand{\footrulewidth}{0pt}
}

\fancypagestyle{plain}{
  \fancyhf{}
  \fancyhead[C]{\textbf{\thepage}\\[-0.2ex]\textbf{\doccode}}
  \fancyfoot[C]{\pagefooter}
  \renewcommand{\headrulewidth}{0pt}
  \renewcommand{\footrulewidth}{0pt}
}

\newcommand{\leftapprovalstamp}[1]{%
  \begingroup
  \renewcommand{\arraystretch}{1.0}
  \setlength{\tabcolsep}{0pt}
  \footnotesize
  \begin{tabular}{|>{\centering\arraybackslash}m{0.65cm}|>{\centering\arraybackslash}m{1.05cm}|}
    \hline
    \rotatebox{90}{Подп. и дата} & \rule{0pt}{1.10cm}\\
    \hline
    \rotatebox{90}{Инв. № дубл.} & \rule{0pt}{1.00cm}\\
    \hline
    \rotatebox{90}{Взам. инв. №} & \rule{0pt}{1.00cm}\\
    \hline
    \rotatebox{90}{Подп. и дата} & \rule{0pt}{1.10cm}\\
    \hline
    \rotatebox{90}{Инв. № подл.} &
    \rotatebox{90}{\scalebox{0.70}{#1}}\rule{0pt}{0.95cm}\\
    \hline
  \end{tabular}
  \endgroup
}

\newcommand{\placestamp}[2]{%
  \begin{textblock*}{19mm}(1mm,#2)
    \leftapprovalstamp{#1}
  \end{textblock*}
}

\newcommand{\registrationblankrow}{%
  \rule{0pt}{0.78cm} & & & & & & & & & \\
  \hline
}

\newcommand{\registrationtable}{%
  \begingroup
  \setlength{\tabcolsep}{2pt}%
  \renewcommand{\arraystretch}{1.0}
  \footnotesize
  \centering
  \begin{tabular}{|>{\centering\arraybackslash}p{0.75cm}|>{\centering\arraybackslash}p{1.55cm}|>{\centering\arraybackslash}p{1.55cm}|>{\centering\arraybackslash}p{1.55cm}|>{\centering\arraybackslash}p{1.7cm}|>{\centering\arraybackslash}p{1.9cm}|>{\centering\arraybackslash}p{2.2cm}|>{\centering\arraybackslash}p{2.4cm}|>{\centering\arraybackslash}p{1.0cm}|>{\centering\arraybackslash}p{1.0cm}|}
    \hline
    \multicolumn{10}{|c|}{\textbf{Лист регистрации изменений}}\\
    \hline
    \multicolumn{5}{|c|}{\makecell[c]{Номера листов\\(страниц)}} & & & & & \\
    \hline
    Изм. & \makecell{Изме-\\ненных} & \makecell{Заме-\\ненных} & Новых & \makecell{Аннулиро-\\ванных} &
    \makecell[c]{Всего\\листов\\(страниц)\\в докум.} &
    \makecell[c]{№\\документа} &
    \makecell[c]{Входящий\\№\\сопроводит.\\документа\\и\\дата} &
    \makecell[c]{Подп.} &
    \makecell[c]{Дата} \\
    \hline
    \registrationblankrow
    \registrationblankrow
    \registrationblankrow
    \registrationblankrow
    \registrationblankrow
    \registrationblankrow
    \registrationblankrow
    \registrationblankrow
    \registrationblankrow
    \registrationblankrow
    \registrationblankrow
    \registrationblankrow
    \registrationblankrow
    \registrationblankrow
  \end{tabular}
  \endgroup
}

\begin{document}
\hypersetup{pageanchor=false}
\selectlanguage{russian}

\thispagestyle{empty}
\placestamp{RU.17701729.05.05-01 ТЗ 01-1}{146mm}
\vspace*{-2.1cm}

\begin{center}
  {\small\textbf{ПРАВИТЕЛЬСТВО РОССИЙСКОЙ ФЕДЕРАЦИИ}}\\
  {\small\textbf{ФЕДЕРАЛЬНОЕ ГОСУДАРСТВЕННОЕ АВТОНОМНОЕ}}\\
  {\small\textbf{ОБРАЗОВАТЕЛЬНОЕ УЧРЕЖДЕНИЕ ВЫСШЕГО ОБРАЗОВАНИЯ}}\\
  {\small\textbf{НАЦИОНАЛЬНЫЙ ИССЛЕДОВАТЕЛЬСКИЙ УНИВЕРСИТЕТ}}\\
  {\small\textbf{«ВЫСШАЯ ШКОЛА ЭКОНОМИКИ»}}\\[0.35cm]
  Факультет компьютерных наук\\
  Образовательная программа «Программная инженерия»
\end{center}

\vspace{0.6cm}
\noindent
\begin{minipage}[t]{0.49\textwidth}
  \centering
  \textbf{СОГЛАСОВАНО}\\
  Руководитель,\\
  Доцент департамента «Программной инженерии»\\[0.55cm]
  \underline{\hspace{2.7cm}} А.Д. Брейман\\
  «\underline{\hspace{0.9cm}}» \underline{\hspace{2.1cm}} 2026 г.
\end{minipage}\hfill
\begin{minipage}[t]{0.49\textwidth}
  \centering
  \textbf{УТВЕРЖДАЮ}\\
  Академический руководитель\\
  образовательной программы\\
  «Программная инженерия»,\\
  старший преподаватель департамента\\
  программной инженерии\\[0.2cm]
  \underline{\hspace{2.7cm}} Н.А. Павлочев\\
  «\underline{\hspace{0.9cm}}» \underline{\hspace{2.1cm}} 2026 г.
\end{minipage}

\vspace{0.75cm}
\begin{center}
  {\large\textbf{\doctitle}}\\[0.22cm]
  {\large\textbf{Техническое задание}}\\[0.18cm]
  {\large\textbf{ЛИСТ УТВЕРЖДЕНИЯ}}\\[0.18cm]
  {\large\textbf{\doccodeapproval}}
\end{center}

\vspace{1.15cm}
\begin{flushright}
  Исполнитель:\\[0.25cm]
  студент группы БПИ223\\[0.35cm]
  \underline{\hspace{2.5cm}} И.И. Новик\\[0.25cm]
  «\underline{\hspace{0.9cm}}» \underline{\hspace{2.8cm}} 2026 г.
\end{flushright}

\begin{textblock*}{\paperwidth}(0mm,281mm)
  \centering{\large\textbf{2026}}
\end{textblock*}
\newpage

\thispagestyle{empty}
\placestamp{RU.17701729.05.05-01 ТЗ 01-1-ЛУ}{144mm}

\vspace*{-0.6cm}
\begin{flushleft}
  \textbf{УТВЕРЖДЕН}\\[0.2cm]
  \textbf{\doccodeapproval}
\end{flushleft}

\vspace{1.4cm}
\begin{center}
  {\large\textbf{\doctitle}}\\[0.22cm]
  {\Large\textbf{Техническое задание}}\\[0.22cm]
  {\Large\textbf{\doccodebase}}\\[0.2cm]
  {\Large\textbf{Листов \pageref{lastnumberedpage}}}
\end{center}

\begin{textblock*}{\paperwidth}(0mm,281mm)
  \centering{\large\textbf{2026}}
\end{textblock*}

\clearpage
\hypersetup{pageanchor=true}
\setcounter{page}{2}
\pagestyle{tzstyle}
\tableofcontents
\clearpage

\section*{АННОТАЦИЯ}
\addcontentsline{toc}{section}{АННОТАЦИЯ}
Техническое задание является основным документом, который определяет цель разработки,
набор требований, структуру и порядок создания программного продукта. По этому документу
проводятся разработка, испытания и приемка результата без двусмысленных трактовок.

Настоящее техническое задание на разработку Андроид-приложения для совместного
планирования поездок содержит следующие разделы:
\begin{enumerate}
  \item «Введение»;
  \item «Основания для разработки»;
  \item «Назначение разработки»;
  \item «Требования к программе»;
  \item «Требования к программной документации»;
  \item «Технико-экономические показатели»;
  \item «Стадии разработки»;
  \item «Порядок контроля и приемки».
\end{enumerate}

Разделы документа раскрывают область применения и назначение приложения, проверяемые
функциональные требования и ограничения, требования к документации и порядок приемки.

\clearpage
\section*{ГЛОССАРИЙ}
\addcontentsline{toc}{section}{ГЛОССАРИЙ}
\textbf{Техническое задание (ТЗ)} --- документ, определяющий требования к разработке, проверке и приемке результата.\\
\textbf{Программа и методика испытаний (ПМИ)} --- документ, определяющий порядок проверки пользовательских сценариев приложения.\\
\textbf{Выпускная квалификационная работа (ВКР)} --- итоговая работа, описывающая результаты проектирования, реализации и проверки приложения.\\
\textbf{Искусственный интеллект (ИИ)} --- серверный модуль генерации текстовых рекомендаций для поездки на основе внешней языковой модели и параметров, заданных пользователем.\\
\textbf{Приложение (CoTrip)} --- мобильное Android-приложение для совместного планирования поездок.\\
\textbf{Пользователь} --- авторизованный человек, выполняющий действия в приложении.\\
\textbf{Поездка} --- сущность с датами, участниками, идеями, маршрутом и расходами.\\
\textbf{Владелец поездки} --- пользователь, создавший поездку.\\
\textbf{Участник поездки} --- пользователь, присоединившийся к поездке.\\
\textbf{Идея} --- предложение активности или места до включения в маршрут.\\
\textbf{Комментарий идеи} --- сообщение в обсуждении конкретной идеи.\\
\textbf{Активность} --- элемент маршрута на конкретный день.\\
\textbf{Расход} --- запись о сумме, плательщике и распределении между участниками.\\
\textbf{Баланс} --- итоговое финансовое состояние пользователя в рамках поездки.\\
\textbf{Дни вне диапазона} --- дни маршрута, которые выходят за текущий диапазон дат поездки.\\
\textbf{Подсказка ИИ} --- предложение активности для поездки, сформированное модулем искусственного интеллекта.\\
\textbf{Офлайн-очередь} --- локально сохраненные изменения, отправляемые после восстановления сети.\\
\textbf{Глубокая ссылка} --- ссылка, открывающая конкретный экран приложения.

\clearpage
\section{Введение}

\subsection{Наименование программы}
\begin{itemize}
  \item Наименование на русском: \textbf{Андроид-приложение для совместного планирования поездок}.
  \item Наименование на английском: \textbf{Android Application for Collaborative Trip Planning}.
  \item Краткое имя: \textbf{CoTrip}.
\end{itemize}

\subsection{Область применения}
Приложение предназначено для групп пользователей, планирующих поездку: идеи, обсуждение,
маршрут, расходы, погода и подсказки ИИ доступны в одном интерфейсе.

\clearpage
\section{Основания для разработки}
\subsection{Документ, на основании которого ведется разработка}
Разработка выполняется в рамках образовательной программы бакалавриата 09.03.04
«Программная инженерия» факультета компьютерных наук НИУ ВШЭ.

Основанием для разработки являются учебный план подготовки бакалавров по направлению
09.03.04 «Программная инженерия» и утвержденная тема выпускной квалификационной работы:
«Андроид-приложение для совместного планирования поездок».

\subsection{Наименование темы разработки}
Наименование программы на русском языке: «Андроид-приложение для совместного планирования поездок».

Наименование программы на английском языке: «Android Application for Collaborative Trip Planning».

\clearpage
\section{Назначение разработки}

\subsection{Функциональное назначение}
Приложение должно поддерживать пользовательские сценарии:
\begin{enumerate}
  \item вход и управление профилем;
  \item создание и ведение поездок;
  \item приглашение и присоединение участников;
  \item работа с идеями и обсуждениями;
  \item построение маршрута по дням;
  \item учет расходов и просмотр баланса;
  \item просмотр погоды;
  \item генерация и сохранение подсказок ИИ;
  \item работа при кратковременной потере сети.
\end{enumerate}

\subsection{Эксплуатационное назначение}
Приложение используется на Android-смартфонах пользователями без специальной технической подготовки.

\clearpage
\section{Требования к программе}

\subsection{Функциональные требования}

\subsubsection{Авторизация и профиль}
\begin{enumerate}
  \item пользователь должен иметь возможность войти через Google;
  \item после успешного входа система должна открыть экран списка поездок;
  \item пользователь должен иметь возможность изменить имя и фото профиля;
  \item система должна сохранять измененные имя и фото и отображать их после повторного открытия настроек;
  \item пользователь должен иметь возможность выйти из аккаунта;
  \item после выхода система должна открыть экран входа;
  \item пользователь должен иметь возможность удалить профиль;
  \item после удаления профиля система должна открыть экран входа;
  \item система не должна сохранять пустое имя профиля;
  \item при сетевой ошибке система должна отображать сообщение об ошибке.
\end{enumerate}

\subsubsection{Поездки}
\begin{enumerate}
  \item пользователь должен иметь возможность создать, отредактировать, удалить поездку и открыть экран ее деталей;
  \item после создания поездка должна отображаться в списке пользователя;
  \item после редактирования поездки система должна сохранять изменения;
  \item после удаления поездка не должна отображаться в списке;
  \item система не должна создавать или сохранять поездку без названия;
  \item система не должна сохранять поездку, если дата окончания раньше даты начала;
  \item при создании поездки система не должна принимать дату начала раньше текущей даты;
  \item при создании поездки система не должна принимать дату начала позже, чем через один год после текущей даты;
  \item система не должна принимать дату окончания позже, чем через шесть месяцев после даты начала;
  \item при редактировании существующей поездки система должна разрешать сохранение исходной пары дат, даже если эта пара выходит за текущие ограничения диапазона.
\end{enumerate}

\subsubsection{Приглашения и присоединение}
\begin{enumerate}
  \item владелец поездки должен иметь возможность сгенерировать ссылку-приглашение;
  \item пользователь должен иметь возможность присоединиться к поездке по корректной ссылке;
  \item после присоединения поездка должна появляться в списке пользователя с актуальным составом участников;
  \item система не должна присоединять пользователя по ссылке или токену некорректного формата;
  \item повторное присоединение одного и того же пользователя не должно создавать дубликат участника.
\end{enumerate}

\subsubsection{Идеи и обсуждения}
\begin{enumerate}
  \item участник поездки должен иметь возможность создать, отредактировать и удалить идею;
  \item пользователь должен иметь возможность отправить комментарий к идее;
  \item комментарий, отправленный одним участником, должен отображаться у других участников без ручного обновления;
  \item пользователь должен иметь возможность изменить статус идеи на одно из значений: \texttt{ожидает}, \texttt{одобрена}, \texttt{отклонена};
  \item пользователь должен иметь возможность добавить идею в выбранный день маршрута;
  \item система не должна сохранять идею без заголовка;
  \item система не должна отправлять пустой комментарий;
  \item в текущей версии пользовательский интерфейс не должен предоставлять удаление комментариев.
\end{enumerate}

\subsubsection{Маршрут}
\begin{enumerate}
  \item пользователь должен иметь возможность выбрать город для дня поездки;
  \item пользователь должен иметь возможность добавить, отредактировать и удалить активность;
  \item пользователь должен иметь возможность менять порядок активностей и переносить их между днями;
  \item система должна отображать маршрут по дням поездки;
  \item система не должна добавлять активность с пустым названием;
  \item измененный порядок активностей должен сохраняться после повторного открытия экрана;
  \item при изменении дат поездки система должна предоставить пользователю выбор способа обработки дней вне диапазона.
\end{enumerate}

\subsubsection{Расходы и баланс}
\begin{enumerate}
  \item пользователь должен иметь возможность создать расход с выбором плательщика и участников;
  \item пользователь должен иметь возможность изменить статус оплаты по участникам расхода;
  \item система должна отображать расход в списке расходов поездки;
  \item система должна отображать суммы «оплачено», «запланировано» и баланс пользователя;
  \item система должна отображать в деталях расхода доли участников и их статусы оплаты;
  \item система не должна сохранять расход без названия;
  \item система не должна сохранять расход без суммы;
  \item система не должна сохранять расход без выбранных участников.
\end{enumerate}

\subsubsection{Погода}
\begin{enumerate}
  \item пользователь должен иметь возможность открыть экран прогноза погоды для города поездки;
  \item пользователь должен иметь возможность обновить прогноз;
  \item система должна отображать прогноз для выбранного города;
  \item система должна отображать покрытие дат прогноза;
  \item при отсутствии выбранного города или недоступности прогноза система должна отображать поясняющее сообщение о причине.
\end{enumerate}

\subsubsection{Подсказки ИИ}
\begin{enumerate}
  \item пользователь должен иметь возможность сгенерировать список подсказок для выбранного города;
  \item пользователь должен иметь возможность сохранить выбранную подсказку в идеи поездки;
  \item система должна отображать список сгенерированных подсказок;
  \item сохраненная подсказка должна отображаться в списке идей поездки;
  \item сохраненная подсказка должна отображаться как сохраненная.
\end{enumerate}

\subsubsection{Настройки уведомлений}
\begin{enumerate}
  \item пользователь должен иметь возможность включать и выключать категории уведомлений;
  \item система должна сохранять состояния переключателей категорий после повторного открытия экрана;
  \item при включенной категории уведомлений, выданном системном разрешении и доступном канале доставки уведомлений открытие уведомления должно переводить пользователя в связанный раздел приложения;
  \item поддерживаемые типы событий для перехода по уведомлению: \textbf{новый комментарий к идее}, \textbf{новая идея}, \textbf{новый расход}, \textbf{изменение статуса оплаты расхода}.
\end{enumerate}

\subsubsection{Офлайн-поведение}
\begin{enumerate}
  \item при отсутствии сети система должна предоставлять доступ к ранее загруженным данным поездки;
  \item действия пользователя, выполненные без сети, должны сохраняться в офлайн-очередь;
  \item после восстановления сети система должна отправить офлайн-очередь и обновить экраны.
\end{enumerate}

\subsubsection{Локализация интерфейса}
\begin{enumerate}
  \item если системный язык устройства --- русский, интерфейс приложения должен отображаться на русском языке;
  \item при любом другом системном языке интерфейс приложения должен отображаться на английском языке.
\end{enumerate}

\subsection{Временные характеристики}
\begin{enumerate}
  \item в операциях ожидания отображается индикатор загрузки;
  \item пользователь получает визуальную обратную связь до получения ответа сервера.
\end{enumerate}

\subsection{Надежность}
\begin{enumerate}
  \item сетевые и серверные ошибки отображаются сообщением в интерфейсе;
  \item некорректный ввод не приводит к аварийному завершению;
  \item после перезапуска приложения доступ к данным пользователя сохраняется в рамках
  активной сессии или кэша.
\end{enumerate}

\subsection{Условия эксплуатации}
\begin{enumerate}
  \item целевая платформа: Android 8.0+ (API 26+);
  \item для онлайн-сценариев требуется доступ в интернет.
\end{enumerate}

\clearpage
\section{Требования к программной документации}
Комплект документации:
\begin{enumerate}
  \item настоящее техническое задание (ТЗ);
  \item программа и методика испытаний (ПМИ);
  \item выпускная квалификационная работа (ВКР);
  \item текст программы.
\end{enumerate}

\clearpage
\section{Технико-экономические показатели}

\subsection{Ориентировочная экономическая эффективность}
В рамках проекта расчет экономической эффективности программного продукта не производился.

\subsection{Предполагаемая годовая потребность}
В рамках проекта расчет экономической эффективности программного продукта не производился.

\subsection{Экономические преимущества разработки по сравнению с лучшими отечественными и зарубежными образцами или аналогами}
Сопоставление с аналогами:
\begin{enumerate}
  \item Splitwise закрывает учет расходов, но не закрывает идеи, обсуждения и маршрут в одном процессе;
  \item Wanderlog закрывает маршрутную часть, но не закрывает полный финансовый и дискуссионный контур;
  \item TravelSpend закрывает учет расходов, но не закрывает маршрут и коллективное обсуждение;
  \item TripIt агрегирует данные поездки, но не закрывает управляемый цикл «идея --- обсуждение --- маршрут»;
  \item Яндекс Путешествия и RUSSPASS ориентированы на подбор и информационный контент, а не на полный цикл совместного планирования с учетом расходов;
  \item разрабатываемое приложение CoTrip объединяет в одном пространстве приглашения участников, идеи и обсуждения, маршрут по дням, расходы и баланс.
\end{enumerate}

Экономические преимущества разработки относительно указанных аналогов:
\begin{enumerate}
  \item единое приложение исключает необходимость параллельного ведения поездки в нескольких сервисах;
  \item единая модель данных уменьшает дублирование ввода и вероятность расхождения версий плана между участниками;
  \item встроенный контур «приглашение --- идея --- обсуждение --- маршрут --- расход» сокращает число ручных согласований внутри группы;
  \item офлайн-очередь с последующей синхронизацией снижает риск потери пользовательских действий при нестабильной сети.
\end{enumerate}

\clearpage
\section{Стадии разработки}
Стадии и этапы разработки представлены в таблице \ref{tab:stages}.

\begin{footnotesize}
\begin{longtable}{|>{\raggedright\arraybackslash}p{0.15\textwidth}|>{\raggedright\arraybackslash}p{0.20\textwidth}|>{\raggedright\arraybackslash}p{0.28\textwidth}|>{\raggedright\arraybackslash}p{0.12\textwidth}|>{\raggedright\arraybackslash}p{0.12\textwidth}|}
  \caption{Стадии и этапы разработки}\label{tab:stages}\\
  \hline
  \textbf{\makecell[l]{Стадии\\разработки}} & \textbf{\makecell[l]{Этапы\\разработки}} & \textbf{\makecell[l]{Содержание\\работ}} & \textbf{Сроки} & \textbf{\makecell[l]{Исполни-\\тели}} \\
  \hline
  \endfirsthead

  \hline
  \textbf{\makecell[l]{Стадии\\разработки}} & \textbf{\makecell[l]{Этапы\\разработки}} & \textbf{\makecell[l]{Содержание\\работ}} & \textbf{Сроки} & \textbf{\makecell[l]{Исполни-\\тели}} \\
  \hline
  \endhead

  \multirow{2}{=}{1. Техническое задание}
  & Подготовительные работы
  & Постановка задачи;\newline
    Сбор необходимых данных;\newline
    Выбор и обоснование критериев эффективности и качества программы.
  & \multirow{2}{=}{15.10.25 -- 15.12.25}
  & \multirow{2}{=}{Новик И.\,И.} \\
  \cline{2-3}
  & Разработка и утверждение технического задания
  & Определение требований к разрабатываемой программе;\newline
    Определение требований к техническим средствам;\newline
    Определение этапов разработки и документации;\newline
    Согласование и утверждение технического задания.
  &  &  \\
  \hline

  \multirow{3}{=}{2. Рабочий проект}
  & Разработка программы
  & Создание макета приложения;\newline
    Разработка backend;\newline
    Разработка Android-клиента;\newline
    Разработка, согласование и утверждение порядка и методики испытаний.
  & \multirow{3}{=}{15.12.25 -- 01.05.26}
  & \multirow{3}{=}{Новик И.\,И.} \\
  \cline{2-3}
  & Разработка программной документации
  & Разработка программных документов в соответствии с требованиями ГОСТ 19.101-77.
  &  &  \\
  \cline{2-3}
  & Испытание программы
  & Разработка, согласование и утверждение порядка и методики испытаний;\newline
    Корректировка программы и программной документации по результатам испытаний.
  &  &  \\
  \hline

  3. Внедрение
  & Подготовка и передача программы
  & Подготовка программы и программной документации для презентации и защиты.
  & 01.05.26 -- 15.06.26
  & Новик И.\,И. \\
  \hline
\end{longtable}
\end{footnotesize}

\noindent
Разработка данного программного продукта должна быть завершена к 1 июня 2026 г.
Исполнитель -- Новик Илья Иванович, БПИ223, 4 курс факультета компьютерных наук НИУ ВШЭ.

\clearpage
\section{Порядок контроля и приемки}

\subsection{Виды испытаний}
Проверка программного продукта, в том числе на соответствие настоящему техническому
заданию, осуществляется исполнителем согласно «Программе и методике испытаний».
Приемка выполненного проекта осуществляется комиссией из преподавателей департамента
программной инженерии в сроки, утвержденные приказом декана факультета компьютерных наук.

\clearpage
\section*{ПРИЛОЖЕНИЕ 1}
\addcontentsline{toc}{section}{ПРИЛОЖЕНИЕ 1}
\begin{center}
  \textbf{СПИСОК ИСПОЛЬЗУЕМОЙ ЛИТЕРАТУРЫ}
\end{center}
\begin{enumerate}
  \item ГОСТ 19.101-77. Единая система программной документации. Виды программ и программных документов.
  \item ГОСТ 19.102-77. Единая система программной документации. Стадии разработки.
  \item ГОСТ 19.103-77. Единая система программной документации. Обозначения программ и программных документов.
  \item ГОСТ 19.104-78. Единая система программной документации. Основные надписи.
  \item ГОСТ 19.105-78. Единая система программной документации. Общие требования к программным документам.
  \item ГОСТ 19.106-78. Единая система программной документации. Требования к программным документам, выполненным печатным способом.
  \item ГОСТ 19.201-78. Единая система программной документации. Техническое задание. Требования к содержанию и оформлению.
  \item ГОСТ 19.301-79. Единая система программной документации. Программа и методика испытаний. Требования к содержанию и оформлению.
  \item ГОСТ 19.401-78. Единая система программной документации. Текст программы. Требования к содержанию и оформлению.
  \item ГОСТ 19.404-79. Единая система программной документации. Пояснительная записка. Требования к содержанию и оформлению.
  \item ГОСТ 19.505-79. Единая система программной документации. Руководство оператора. Требования к содержанию и оформлению.
  \item ГОСТ 19.603-78. Единая система программной документации. Общие правила внесения изменений.
  \item ГОСТ 19.604-78. Единая система программной документации. Правила внесения изменений в программные документы, выполненные печатным способом.
  \item Splitwise. URL: \url{https://www.splitwise.com/} (дата обращения: 19.03.2026).
  \item Wanderlog. URL: \url{https://wanderlog.com/} (дата обращения: 19.03.2026).
  \item TravelSpend. URL: \url{https://travel-spend.com/} (дата обращения: 19.03.2026).
  \item TripIt. URL: \url{https://www.tripit.com/} (дата обращения: 19.03.2026).
  \item Яндекс. «Яндекс Путешествия упростили совместные поездки», 03.04.2025. URL: \url{https://yandex.ru/company/news/01-03-04-2025} (дата обращения: 19.03.2026).
  \item RUSSPASS (Google Play). URL: \url{https://play.google.com/store/apps/details?id=ru.russpass.tourist} (дата обращения: 19.03.2026).
  \item Android Developers. Build your first app. URL: \url{https://developer.android.com/training/basics/firstapp} (дата обращения: 19.03.2026).
  \item Ktor Documentation. Server development. URL: \url{https://ktor.io/docs/server-create-a-new-project.html} (дата обращения: 19.03.2026).
  \item JetBrains. Kotlin Documentation. URL: \url{https://kotlinlang.org/docs/home.html} (дата обращения: 19.03.2026).
  \item Android Developers. Jetpack Compose. URL: \url{https://developer.android.com/compose} (дата обращения: 19.03.2026).
  \item Android Developers. Dependency Injection with Hilt. URL: \url{https://developer.android.com/training/dependency-injection/hilt-android} (дата обращения: 19.03.2026).
  \item Square. Retrofit. URL: \url{https://square.github.io/retrofit/} (дата обращения: 19.03.2026).
  \item OpenWeather. Weather API. URL: \url{https://openweathermap.org/api} (дата обращения: 19.03.2026).
  \item Yandex Cloud. YandexGPT (Foundation Models). URL: \url{https://yandex.cloud/en/docs/ai-studio/concepts/generation/models} (дата обращения: 19.03.2026).
  \item Android Developers. WorkManager. URL: \url{https://developer.android.com/topic/libraries/architecture/workmanager} (дата обращения: 19.03.2026).
\end{enumerate}

\label{lastnumberedpage}

\clearpage
\pagestyle{empty}
\thispagestyle{empty}
\pagenumbering{gobble}
\vspace*{-1.1cm}
\begin{center}
  \textbf{\Large ЛИСТ РЕГИСТРАЦИИ ИЗМЕНЕНИЙ}
\end{center}
\vspace{0.35cm}
\registrationtable

\end{document}
